% analysis/design_principles.tex

\section{Design Principles of Alignment-Based Information Modelling}

% ============================================================
\subsection*{Motivation}
% ============================================================

\begin{remark}
The AIM framework integrates geometry, coordinate systems, realization,
and optimization into a unified model.
\end{remark}

\begin{remark}
From this integration, a set of fundamental design principles emerges.
\end{remark}

---

% ============================================================
\subsection{Principle 1: Curvature-Centric Modelling}
% ============================================================

\begin{proposition}
Railway alignments should be modelled primarily through curvature
functions \(\kappa(s)\), not through point-based geometry.
\end{proposition}

\begin{remark}
Curvature provides:
\begin{itemize}
\item intrinsic representation,
\item direct link to dynamics,
\item natural integration with optimization.
\end{itemize}
\end{remark}

---

% ============================================================
\subsection{Principle 2: Separation of Representation and Realization}
% ============================================================

\begin{proposition}
The theoretical \defterm{alignment} model must be separated from its physical
realization.
\end{proposition}

\begin{remark}
The realization operator introduces:
\begin{itemize}
\item smoothing,
\item constraints,
\item deviations from the ideal model.
\end{itemize}
\end{remark}

\begin{remark}
Ignoring this separation leads to systematic modelling errors.
\end{remark}

---

% ============================================================
\subsection{Principle 3: Hybrid Modelling}
% ============================================================

\begin{proposition}
Alignment models should combine parametric structure with local
correction terms.
\end{proposition}

\begin{remark}
This allows:
\begin{itemize}
\item interpretable design intent,
\item robust handling of imperfect data,
\item flexible refinement.
\end{itemize}
\end{remark}

---

% ============================================================
\subsection{Principle 4: Operator-Based Architecture}
% ============================================================

\begin{proposition}
Alignment modelling should be formulated as a composition of operators:
\[
\mathcal{P}
=
\mathcal{W}^{-1}
\circ
\mathcal{R}
\circ
\mathcal{M}
\circ
\mathcal{W}
\circ
\mathcal{C}.
\]
\end{proposition}

\begin{remark}
Each operator represents a distinct domain:
\begin{itemize}
\item coordinate systems,
\item geometric modelling,
\item physical processes,
\item data transformation.
\end{itemize}
\end{remark}

---

% ============================================================
\subsection{Principle 5: Design in Realized Space}
% ============================================================

\begin{proposition}
Optimization objectives must be evaluated on realized geometry, not on
theoretical models.
\end{proposition}

\begin{remark}
The relevant mapping is:
\[
\kappa_{\mathrm{target}}
\;\mapsto\;
\gamma_{\mathrm{real}}.
\]
\end{remark}

\begin{remark}
Design must account for the transformation induced by realization.
\end{remark}

---

% ============================================================
\subsection{Principle 6: Exploitation of Structure}
% ============================================================

\begin{proposition}
The intrinsic block and sparsity structure of the problem must be
exploited for efficient computation.
\end{proposition}

\begin{remark}
This includes:
\begin{itemize}
\item locality in arc-length,
\item block decomposition (curvature vs. cant),
\item sparse Jacobians.
\end{itemize}
\end{remark}

---

% ============================================================
\subsection{Principle 7: Scale Awareness}
% ============================================================

\begin{proposition}
Different components of the pipeline must be evaluated at appropriate
scales.
\end{proposition}

\begin{remark}
\begin{itemize}
\item CRS dominates at large scale,
\item parametric modelling at medium scale,
\item corrections and realization at small scale.
\end{itemize}
\end{remark}

---

% ============================================================
\subsection{Principle 8: Measurement Integration}
% ============================================================

\begin{proposition}
Measured data must be integrated through the world-to-track
transformation, not directly in world space.
\end{proposition}

\begin{remark}
This ensures consistency with the intrinsic alignment representation.
\end{remark}

---

% ============================================================
\subsection{Principle 9: Regularization as Necessity}
% ============================================================

\begin{proposition}
Regularization is an essential component of alignment modelling and
optimization.
\end{proposition}

\begin{remark}
It stabilizes:
\begin{itemize}
\item inverse problems,
\item hybrid corrections,
\item numerical optimization.
\end{itemize}
\end{remark}

---

% ============================================================
\subsection{Principle 10: Acceptance of Non-Uniqueness}
% ============================================================

\begin{proposition}
Alignment solutions are not unique and must be evaluated based on
engineering criteria.
\end{proposition}

\begin{remark}
Different models may produce nearly identical realized geometry.
\end{remark}

\begin{remark}
Thus, solution quality must be defined through:
\begin{itemize}
\item dynamics,
\item constraints,
\item robustness.
\end{itemize}
\end{remark}

---

% ============================================================
\subsection{Summary}
% ============================================================

\begin{summary}
The AIM framework is governed by a set of design principles that
emerge from the interaction of geometry, physics, data, and computation.

These principles define:
\begin{itemize}
\item how alignments should be represented,
\item how they should be transformed,
\item how they should be optimized,
\item and how their limitations must be handled.
\end{itemize}

Together, they provide a coherent foundation for alignment-based
information modelling as a unified engineering and computational
discipline.
\end{summary}
